% ============================================================
%  mycommand.tex —— 自定义命令
%  在 template.tex 中通过 % ============================================================
%  mycommand.tex —— 自定义命令
%  在 template.tex 中通过 % ============================================================
%  mycommand.tex —— 自定义命令
%  在 template.tex 中通过 % ============================================================
%  mycommand.tex —— 自定义命令
%  在 template.tex 中通过 \input{mycommand.tex} 引入
% ============================================================

% ---- 高亮强调 ----
\newcommand{\highlightit}[1]{%
	\colorbox{yellow!60}{\textcolor{black}{#1}}%
}

% ---- 粗体向量（标量） ----
\newcommand{\bsa}{\boldsymbol{a}}
\newcommand{\bsb}{\boldsymbol{b}}
\newcommand{\bse}{\boldsymbol{e}}
\newcommand{\bsf}{\boldsymbol{f}}
\newcommand{\bsg}{\boldsymbol{g}}
\newcommand{\bsn}{\boldsymbol{n}}
\newcommand{\bsq}{\boldsymbol{q}}
\newcommand{\bssf}{\boldsymbol{f}}
\newcommand{\bst}{\boldsymbol{t}}
\newcommand{\bsu}{\boldsymbol{u}}
\newcommand{\bsv}{\boldsymbol{v}}
\newcommand{\bsw}{\boldsymbol{w}}
\newcommand{\bsx}{\boldsymbol{x}}

% ---- 粗体向量（大写/物理量） ----
\newcommand{\bsA}{\boldsymbol{A}}
\newcommand{\bsB}{\boldsymbol{B}}
\newcommand{\bsD}{\boldsymbol{D}}
\newcommand{\bsE}{\boldsymbol{E}}
\newcommand{\bsF}{\boldsymbol{F}}
\newcommand{\bsH}{\boldsymbol{H}}
\newcommand{\bsJ}{\boldsymbol{J}}
\newcommand{\bsM}{\boldsymbol{M}}
\newcommand{\bsV}{\boldsymbol{V}}

% ---- 微分算子 ----
\newcommand{\bcurl}{\boldsymbol{\mathrm{curl}}\,}
\newcommand{\diver}{\mathrm{div}\,}
\newcommand{\rot}{\mathrm{rot}\,}

% ---- 通用粗体快捷命令 ----
\newcommand{\bs}[1]{\boldsymbol{#1}}

% ---- 个人简介页辅助命令 ----

% 深色圆形序号
\newcommand{\circnum}[1]{%
	\tikz[baseline=(n.base)]{%
		\node[circle, fill=structure.fg, inner sep=1.8pt,
		minimum size=16pt,
		font=\fontsize{6.5}{6.5}\selectfont\bfseries\color{white}]
		(n) {#1};}%
}

% 研究方向标签（行内带圆角边框）
\newcommand{\researchtag}[1]{%
	\tcbox[on line, arc=4pt,
	colback=white, colframe=tagborder,
	boxrule=0.8pt, top=2pt, bottom=2pt, left=6pt, right=6pt,
	fontupper=\small\bfseries\sffamily\color{structure.fg}]{#1}%
}

% 研究主线条目内容框（行内浅色背景）
\newcommand{\resbox}[2]{%  #1=宽度, #2=文字
	\tcbox[on line, arc=2pt,
	colback=panelbg!50, colframe=panelbg,
	boxrule=0.5pt, top=3pt, bottom=3pt, left=5pt, right=5pt,
	width=#1, fontupper=\fontsize{7.5}{9.5}\selectfont]{#2}%
} 引入
% ============================================================

% ---- 高亮强调 ----
\newcommand{\highlightit}[1]{%
	\colorbox{yellow!60}{\textcolor{black}{#1}}%
}

% ---- 粗体向量（标量） ----
\newcommand{\bsa}{\boldsymbol{a}}
\newcommand{\bsb}{\boldsymbol{b}}
\newcommand{\bse}{\boldsymbol{e}}
\newcommand{\bsf}{\boldsymbol{f}}
\newcommand{\bsg}{\boldsymbol{g}}
\newcommand{\bsn}{\boldsymbol{n}}
\newcommand{\bsq}{\boldsymbol{q}}
\newcommand{\bssf}{\boldsymbol{f}}
\newcommand{\bst}{\boldsymbol{t}}
\newcommand{\bsu}{\boldsymbol{u}}
\newcommand{\bsv}{\boldsymbol{v}}
\newcommand{\bsw}{\boldsymbol{w}}
\newcommand{\bsx}{\boldsymbol{x}}

% ---- 粗体向量（大写/物理量） ----
\newcommand{\bsA}{\boldsymbol{A}}
\newcommand{\bsB}{\boldsymbol{B}}
\newcommand{\bsD}{\boldsymbol{D}}
\newcommand{\bsE}{\boldsymbol{E}}
\newcommand{\bsF}{\boldsymbol{F}}
\newcommand{\bsH}{\boldsymbol{H}}
\newcommand{\bsJ}{\boldsymbol{J}}
\newcommand{\bsM}{\boldsymbol{M}}
\newcommand{\bsV}{\boldsymbol{V}}

% ---- 微分算子 ----
\newcommand{\bcurl}{\boldsymbol{\mathrm{curl}}\,}
\newcommand{\diver}{\mathrm{div}\,}
\newcommand{\rot}{\mathrm{rot}\,}

% ---- 通用粗体快捷命令 ----
\newcommand{\bs}[1]{\boldsymbol{#1}}

% ---- 个人简介页辅助命令 ----

% 深色圆形序号
\newcommand{\circnum}[1]{%
	\tikz[baseline=(n.base)]{%
		\node[circle, fill=structure.fg, inner sep=1.8pt,
		minimum size=16pt,
		font=\fontsize{6.5}{6.5}\selectfont\bfseries\color{white}]
		(n) {#1};}%
}

% 研究方向标签（行内带圆角边框）
\newcommand{\researchtag}[1]{%
	\tcbox[on line, arc=4pt,
	colback=white, colframe=tagborder,
	boxrule=0.8pt, top=2pt, bottom=2pt, left=6pt, right=6pt,
	fontupper=\small\bfseries\sffamily\color{structure.fg}]{#1}%
}

% 研究主线条目内容框（行内浅色背景）
\newcommand{\resbox}[2]{%  #1=宽度, #2=文字
	\tcbox[on line, arc=2pt,
	colback=panelbg!50, colframe=panelbg,
	boxrule=0.5pt, top=3pt, bottom=3pt, left=5pt, right=5pt,
	width=#1, fontupper=\fontsize{7.5}{9.5}\selectfont]{#2}%
} 引入
% ============================================================

% ---- 高亮强调 ----
\newcommand{\highlightit}[1]{%
	\colorbox{yellow!60}{\textcolor{black}{#1}}%
}

% ---- 粗体向量（标量） ----
\newcommand{\bsa}{\boldsymbol{a}}
\newcommand{\bsb}{\boldsymbol{b}}
\newcommand{\bse}{\boldsymbol{e}}
\newcommand{\bsf}{\boldsymbol{f}}
\newcommand{\bsg}{\boldsymbol{g}}
\newcommand{\bsn}{\boldsymbol{n}}
\newcommand{\bsq}{\boldsymbol{q}}
\newcommand{\bssf}{\boldsymbol{f}}
\newcommand{\bst}{\boldsymbol{t}}
\newcommand{\bsu}{\boldsymbol{u}}
\newcommand{\bsv}{\boldsymbol{v}}
\newcommand{\bsw}{\boldsymbol{w}}
\newcommand{\bsx}{\boldsymbol{x}}

% ---- 粗体向量（大写/物理量） ----
\newcommand{\bsA}{\boldsymbol{A}}
\newcommand{\bsB}{\boldsymbol{B}}
\newcommand{\bsD}{\boldsymbol{D}}
\newcommand{\bsE}{\boldsymbol{E}}
\newcommand{\bsF}{\boldsymbol{F}}
\newcommand{\bsH}{\boldsymbol{H}}
\newcommand{\bsJ}{\boldsymbol{J}}
\newcommand{\bsM}{\boldsymbol{M}}
\newcommand{\bsV}{\boldsymbol{V}}

% ---- 微分算子 ----
\newcommand{\bcurl}{\boldsymbol{\mathrm{curl}}\,}
\newcommand{\diver}{\mathrm{div}\,}
\newcommand{\rot}{\mathrm{rot}\,}

% ---- 通用粗体快捷命令 ----
\newcommand{\bs}[1]{\boldsymbol{#1}}

% ---- 个人简介页辅助命令 ----

% 深色圆形序号
\newcommand{\circnum}[1]{%
	\tikz[baseline=(n.base)]{%
		\node[circle, fill=structure.fg, inner sep=1.8pt,
		minimum size=16pt,
		font=\fontsize{6.5}{6.5}\selectfont\bfseries\color{white}]
		(n) {#1};}%
}

% 研究方向标签（行内带圆角边框）
\newcommand{\researchtag}[1]{%
	\tcbox[on line, arc=4pt,
	colback=white, colframe=tagborder,
	boxrule=0.8pt, top=2pt, bottom=2pt, left=6pt, right=6pt,
	fontupper=\small\bfseries\sffamily\color{structure.fg}]{#1}%
}

% 研究主线条目内容框（行内浅色背景）
\newcommand{\resbox}[2]{%  #1=宽度, #2=文字
	\tcbox[on line, arc=2pt,
	colback=panelbg!50, colframe=panelbg,
	boxrule=0.5pt, top=3pt, bottom=3pt, left=5pt, right=5pt,
	width=#1, fontupper=\fontsize{7.5}{9.5}\selectfont]{#2}%
} 引入
% ============================================================

% ---- 高亮强调 ----
\newcommand{\highlightit}[1]{%
	\colorbox{yellow!60}{\textcolor{black}{#1}}%
}

% ---- 粗体向量（标量） ----
\newcommand{\bsa}{\boldsymbol{a}}
\newcommand{\bsb}{\boldsymbol{b}}
\newcommand{\bse}{\boldsymbol{e}}
\newcommand{\bsf}{\boldsymbol{f}}
\newcommand{\bsg}{\boldsymbol{g}}
\newcommand{\bsn}{\boldsymbol{n}}
\newcommand{\bsq}{\boldsymbol{q}}
\newcommand{\bssf}{\boldsymbol{f}}
\newcommand{\bst}{\boldsymbol{t}}
\newcommand{\bsu}{\boldsymbol{u}}
\newcommand{\bsv}{\boldsymbol{v}}
\newcommand{\bsw}{\boldsymbol{w}}
\newcommand{\bsx}{\boldsymbol{x}}

% ---- 粗体向量（大写/物理量） ----
\newcommand{\bsA}{\boldsymbol{A}}
\newcommand{\bsB}{\boldsymbol{B}}
\newcommand{\bsD}{\boldsymbol{D}}
\newcommand{\bsE}{\boldsymbol{E}}
\newcommand{\bsF}{\boldsymbol{F}}
\newcommand{\bsH}{\boldsymbol{H}}
\newcommand{\bsJ}{\boldsymbol{J}}
\newcommand{\bsM}{\boldsymbol{M}}
\newcommand{\bsV}{\boldsymbol{V}}

% ---- 微分算子 ----
\newcommand{\bcurl}{\boldsymbol{\mathrm{curl}}\,}
\newcommand{\diver}{\mathrm{div}\,}
\newcommand{\rot}{\mathrm{rot}\,}

% ---- 通用粗体快捷命令 ----
\newcommand{\bs}[1]{\boldsymbol{#1}}

% ---- 个人简介页辅助命令 ----

% 深色圆形序号
\newcommand{\circnum}[1]{%
	\tikz[baseline=(n.base)]{%
		\node[circle, fill=structure.fg, inner sep=1.8pt,
		minimum size=16pt,
		font=\fontsize{6.5}{6.5}\selectfont\bfseries\color{white}]
		(n) {#1};}%
}

% 研究方向标签（行内带圆角边框）
\newcommand{\researchtag}[1]{%
	\tcbox[on line, arc=4pt,
	colback=white, colframe=tagborder,
	boxrule=0.8pt, top=2pt, bottom=2pt, left=6pt, right=6pt,
	fontupper=\small\bfseries\sffamily\color{structure.fg}]{#1}%
}

% 研究主线条目内容框（行内浅色背景）
\newcommand{\resbox}[2]{%  #1=宽度, #2=文字
	\tcbox[on line, arc=2pt,
	colback=panelbg!50, colframe=panelbg,
	boxrule=0.5pt, top=3pt, bottom=3pt, left=5pt, right=5pt,
	width=#1, fontupper=\fontsize{7.5}{9.5}\selectfont]{#2}%
}